\chapter{Background and Related Work}
\label{ch:background}

This chapter positions \polyforge{} against four bodies of work: autoscaling
for inference serving, semantic caching, multi-tenant fairness, and the
production LLM-serving stack of 2025--26. One reading rule applies
throughout, and it is the same rule the evaluation enforces
(Chapter~\ref{ch:methodology}): \textbf{numbers from other papers are their
numbers on their substrate} --- real GPU fleets, their traces, their
baselines --- while \polyforge{}'s numbers are decision quality on its own
matrix and traces. No comparison in this chapter is a head-to-head cell; the
comparison is \emph{shape against shape}. Where a published system is
reproducible at simulator scale, \polyforge{} carries it as a tuned in-matrix
baseline instead (HPA, KEDA, FIRM-replica, GPTCache-LRU, static, and iso-cost
variants).

\section{Autoscaling for inference serving}

The Kubernetes Horizontal Pod Autoscaler~\citep{k8shpa} and
KEDA~\citep{keda} are the production defaults: reactive controllers on
utilization or event-lag signals. FIRM~\citep{firm2020} added fine-grained,
SLO-oriented resource management for microservices.

The LLM-serving generation is forecast-aware and hierarchical.
SageServe~\citep{sageserve2025} reports ${\sim}25\%$ GPU-hour savings on 10M
production requests using demand forecasting; Chiron~\citep{chiron2025}
reports up to $+90\%$ SLO attainment and $+70\%$ GPU efficiency via
hierarchical autoscaling; Aladdin~\citep{aladdin2024} jointly places and
scales for up to $-71\%$ serving cost; an SLO-driven Kubernetes autoscaler
\citep{k8sslo2025} reports $-31\%$ violation duration at $-18\%$ cost.
Joint horizontal-plus-vertical scaling~\citep{taleoftwoscales2024} and
SLO-constrained joint allocation~\citep{jointalloc2026} widen a single
layer's action space; INFaaS~\citep{infaas2019} pioneered model-variant
selection (``model-less'' serving); Faro~\citep{faro2024} targets on-premise
inference clusters.

Two observations motivate this thesis. First, the closest systems to
``joint'' --- Aladdin (placement$+$scaling), the horizontal$+$vertical
scalers, INFaaS-style variant selection --- each co-optimize \emph{within}
one layer. None of them controls a semantic cache, and none carries
cross-tenant fairness in its objective. Second, the published wins are
typically measured against default or non-SLO-tuned baselines; the bar in
this thesis is deliberately higher --- every baseline is grid-tuned on the
paper's own objective (Section~\ref{sec:method-baselines}).

\section{Semantic caching for LLM serving}

GPTCache~\citep{gptcache2023} established the pattern --- embed the prompt,
serve a stored response on similarity --- reporting 61.6--68.8\% hit rates on
its own benchmark. SCALM~\citep{scalm2024} reports $+63\%$ hits over
GPTCache; InstCache~\citep{instcache2024} pre-populates predicted
instructions and reports 51.34\% on LMSYS-Chat-1M;
MeanCache~\citep{meancache2024} moves the cache user-side and federated;
\citet{semcacheadapt2025} study offline-to-online adaptation of similarity
thresholds.

In all of these, \emph{the cache is the system}: its size is an input, not a
decision variable, and its interaction with capacity and model choice is out
of scope. \polyforge{} treats cache budget as one of three jointly planned
knobs, prices its benefit per tenant in dollars, and adds two measurements
the cache literature does not publish: hit \emph{precision} against the
actually received response on the same real dataset
(Section~\ref{sec:eval-cache}), and a timing side-channel analysis with a
defense frontier (Section~\ref{sec:eval-security};
cf.~the prompt-cache audit of \citet{promptcacheaudit2025}).

\section{Multi-tenant fairness}

VTC~\citep{vtc2024} defined service fairness for LLM serving at token
granularity inside a continuous-batching engine, with a proven $2\times$
service bound; Equinox~\citep{equinox2025}, D\textsuperscript{2}LPM
\citep{d2lpm2025}, and Justitia~\citep{justitia2025} extend it (worst-case
service gaps $-42\%$ vs.\ VTC; locality; efficiency). \polyforge{}'s fairness
sits at a different altitude --- capacity allocation across tenants at the
control-plane level, scored by Jain's index~\citep{jain1984} over per-tenant
SLO satisfaction --- and the claims are scoped to that altitude
(Section~\ref{sec:disc-fairness-altitude}). The evaluation carries a
replica-level transplant of VTC's least-service-first idea as a tuned
baseline (\evfile{vtc\_replica}); beating it on Jain says \polyforge{}
allocates capacity more evenly, not that it schedules tokens more fairly.
Token-level fair schedulers are complementary: they would run \emph{inside}
each replica pool \polyforge{} sizes.

\section{The 2026 serving substrate, and where \polyforge{} sits}
\label{sec:bg-substrate}

The production stack of 2025--26 --- llm-d~\citep{llmd2025},
AIBrix~\citep{aibrix2025}, NVIDIA Dynamo~\citep{dynamo2025}, and the
Kubernetes Gateway API Inference Extension~\citep{gatewayapi2025} ---
actuates on signals \polyforge{}'s system model does not name: queue depth,
in-flight requests, KV-cache pressure, prefill/decode pool split. That is an
altitude difference, not a contradiction. \polyforge{} is the
\emph{portfolio} controller above that actuation layer, and each knob has a
direct target in it:

\begin{table}[htbp]
\centering\small
\begin{tabular}{@{}lll@{}}
\toprule
\textbf{\polyforge{} knob} & \textbf{In the system model} & \textbf{2026 actuation target} \\
\midrule
replicas   & per-tenant capacity      & decode-pool size / vLLM pod count \\
cache MB   & semantic-cache budget    & gateway semantic-cache budget \\
tier       & model class (1:10:100)   & model routing (INFaaS-style; AIBrix) \\
\bottomrule
\end{tabular}
\caption{Substrate mapping of \polyforge{}'s three knobs onto the 2026
LLM-serving stack (from \evfile{docs/RELATED\_WORK.md}~\S4).}
\label{tab:substrate-mapping}
\end{table}

Two design decisions keep the mapping honest: the demand signal sits behind
an interface (\evfile{planner.DemandSource}), so a token-level source (TTFT,
queue depth, KV pressure) replaces the request-rate source without touching
the planner; and latency claims stay request-level, scoped away from
continuous-batching token dynamics. What the 2026 stack does \emph{not} do
--- and where \polyforge{} sits --- is decide jointly, across a tenant
portfolio, in dollars, with per-tenant budget guardrails.

\section{Capability gap summary}

Across all surveyed systems, none co-optimizes replicas, semantic cache, and
model tier against one multi-tenant objective; none accounts per-tenant
dollars with budget guardrails; none publishes a security evaluation of its
cache layer; and none pre-registers its hypotheses or publishes its nulls.
Those four rows --- one architectural, one economic, one adversarial, one
methodological --- are the thesis's position in the literature. The full
capability matrix, including the rows where surveyed systems dominate
\polyforge{} (live GPU fleets, production scale, proven scheduling bounds),
is maintained in \evfile{docs/RELATED\_WORK.md} and summarized honestly in
Section~\ref{sec:disc-limits}.
